% GENERATED FILE. Edit canonical lessons, then run npm run generate:books.
% canonical-source-hash: fnv1a64:2317e63bf1cf0d37
% canonical-lessons: MR-C215-kheladu, MR-C215-prasiksak, MR-C215-cahta, MR-C215-agaman, MR-C215-prasthan

\chapter{Player, Coach, Supporter, Arrival, Departure}
\label{ch:mr-player-coach-supporter-arrival-departure}
\hlchaptermodality{215}

\begin{chapteropening}
\textbf{By the end of this chapter:} \emph{I can say a player, a coach, a supporter, arrival and departure.}

Complete the last lesson of chapter 215: I can say a player, a coach, a supporter, arrival and departure.
\end{chapteropening}

\section[\texorpdfstring{\mr{खेळाडू} (kheḷāḍū)}{खेळाडू (kheḷāḍū)}]{\mr{खेळाडू} (kheḷāḍū) — a player}

\label{lesson:MR-C215-kheladu}



\begin{quote}
Before the new one: say the Marathi for wool, then the Marathi for silk.
\end{quote}

\subsection*{You'll want to know: \mr{खेळाडू}}
\textbf{\mr{खेळाडू}} — \emph{kheḷāḍū} — \textquotedblleft{}a player\textquotedblright{}.

\begin{cousinweb}
\textbf{\mr{खेळाडू}} (\emph{kheḷāḍū}) is \textbf{\mr{खेळ}}, \textquotedblleft{}play\textquotedblright{}, with the ending \emph{-āḍū}, \textquotedblleft{}one who\textquotedblright{}. CDIAL lists Marathi \emph{kheḷ} under a reconstructed \emph{*khēl-}, \textquotedblleft{}play\textquotedblright{} (CDIAL 3918).
\end{cousinweb}

\begin{grammarlens}[title={What you've built}]
One more word for this chapter.
\end{grammarlens}

\subsection*{Your turn}
\begin{itemize}
\raggedright
  \item \emph{Say it:} \emph{kheḷāḍū}
  \item \emph{Say it:} \emph{kheḷāḍū}, once more
  \item \emph{Say it:} say \emph{kheḷāḍū}
  \item \emph{Recall:} say the Marathi for wool, then the Marathi for silk, then say \emph{kheḷāḍū} again
\end{itemize}

\begin{tcolorbox}[breakable,colback=teal!4,colframe=teal!35!black,title={Before you move on}]
What does \mr{खेळाडू} mean? (\textquotedblleft{}A player\textquotedblright{}.) Say it once more.
\end{tcolorbox}
\section[\texorpdfstring{\mr{प्रशिक्षक} (praśikṣak)}{प्रशिक्षक (praśikṣak)}]{\mr{प्रशिक्षक} (praśikṣak) — a coach, a trainer}

\label{lesson:MR-C215-prasiksak}



\begin{quote}
Before the new one: say the Marathi for silk, then the Marathi for a player.
\end{quote}

\subsection*{You'll want to know: \mr{प्रशिक्षक}}
\textbf{\mr{प्रशिक्षक}} — \emph{praśikṣak} — \textquotedblleft{}a coach, a trainer\textquotedblright{}.

\begin{cousinweb}
\textbf{\mr{प्रशिक्षक}} (\emph{praśikṣak}) is a modern word of Sanskrit parts: \emph{pra-}, \textquotedblleft{}forth\textquotedblright{}, and \emph{śikṣaka}, \textquotedblleft{}one who teaches\textquotedblright{}.
\end{cousinweb}

\begin{grammarlens}[title={What you've built}]
One more word for this chapter.
\end{grammarlens}

\subsection*{Your turn}
\begin{itemize}
\raggedright
  \item \emph{Say it:} \emph{praśikṣak}
  \item \emph{Say it:} \emph{praśikṣak}, once more
  \item \emph{Say it:} say \emph{praśikṣak}
  \item \emph{Recall:} say the Marathi for silk, then the Marathi for a player, then say \emph{praśikṣak} again
\end{itemize}

\begin{tcolorbox}[breakable,colback=teal!4,colframe=teal!35!black,title={Before you move on}]
What does \mr{प्रशिक्षक} mean? (\textquotedblleft{}A coach, a trainer\textquotedblright{}.) Say it once more.
\end{tcolorbox}
\section[\texorpdfstring{\mr{चाहता} (cāhtā)}{चाहता (cāhtā)}]{\mr{चाहता} (cāhtā) — a supporter, a fan}

\label{lesson:MR-C215-cahta}



\begin{quote}
Before the new one: say the Marathi for a player, then the Marathi for a coach.
\end{quote}

\subsection*{You'll want to know: \mr{चाहता}}
\textbf{\mr{चाहता}} — \emph{cāhtā} — \textquotedblleft{}a supporter, a fan\textquotedblright{}.

\begin{cousinweb}
\textbf{\mr{चाहता}} (\emph{cāhtā}) is \textquotedblleft{}one who likes\textquotedblright{}, from Marathi \emph{cāhṇẽ}, \textquotedblleft{}to love, like\textquotedblright{}, which CDIAL lists under a reconstructed \emph{*cāh-}, \textquotedblleft{}see, look for, desire\textquotedblright{}, whose source it calls uncertain (CDIAL 4775).
\end{cousinweb}

\begin{grammarlens}[title={What you've built}]
One more word for this chapter.
\end{grammarlens}

\subsection*{Your turn}
\begin{itemize}
\raggedright
  \item \emph{Say it:} \emph{cāhtā}
  \item \emph{Say it:} \emph{cāhtā}, once more
  \item \emph{Say it:} say \emph{cāhtā}
  \item \emph{Recall:} say the Marathi for a player, then the Marathi for a coach, then say \emph{cāhtā} again
\end{itemize}

\begin{tcolorbox}[breakable,colback=teal!4,colframe=teal!35!black,title={Before you move on}]
What does \mr{चाहता} mean? (\textquotedblleft{}A supporter, a fan\textquotedblright{}.) Say it once more.
\end{tcolorbox}
\section[\texorpdfstring{\mr{आगमन} (āgaman)}{आगमन (āgaman)}]{\mr{आगमन} (āgaman) — arrival}

\label{lesson:MR-C215-agaman}



\begin{quote}
Before the new one: say the Marathi for a coach, then the Marathi for a supporter.
\end{quote}

\subsection*{You'll want to know: \mr{आगमन}}
\textbf{\mr{आगमन}} — \emph{āgaman} — \textquotedblleft{}arrival\textquotedblright{}.

\begin{cousinweb}
\textbf{\mr{आगमन}} (\emph{āgaman}) is Sanskrit \emph{āgamana}, \textquotedblleft{}coming, arrival\textquotedblright{}, a tatsama: \emph{ā}, \textquotedblleft{}toward\textquotedblright{}, and \emph{gam}, \textquotedblleft{}to go\textquotedblright{}. CDIAL has no entry for it.
\end{cousinweb}

\begin{grammarlens}[title={What you've built}]
One more word for this chapter.
\end{grammarlens}

\subsection*{Your turn}
\begin{itemize}
\raggedright
  \item \emph{Say it:} \emph{āgaman}
  \item \emph{Say it:} \emph{āgaman}, once more
  \item \emph{Say it:} say \emph{āgaman}
  \item \emph{Recall:} say the Marathi for a coach, then the Marathi for a supporter, then say \emph{āgaman} again
\end{itemize}

\begin{tcolorbox}[breakable,colback=teal!4,colframe=teal!35!black,title={Before you move on}]
What does \mr{आगमन} mean? (\textquotedblleft{}Arrival\textquotedblright{}.) Say it once more.
\end{tcolorbox}
\section[\texorpdfstring{\mr{प्रस्थान} (prasthān)}{प्रस्थान (prasthān)}]{\mr{प्रस्थान} (prasthān) — departure}

\label{lesson:MR-C215-prasthan}



\begin{quote}
Before the new one: say the Marathi for a supporter, then the Marathi for arrival.
\end{quote}

\subsection*{You'll want to know: \mr{प्रस्थान}}
\textbf{\mr{प्रस्थान}} — \emph{prasthān} — \textquotedblleft{}departure\textquotedblright{}.

\begin{cousinweb}
\textbf{\mr{प्रस्थान}} (\emph{prasthān}) is Sanskrit \emph{prasthāna}, \textquotedblleft{}setting out, departure\textquotedblright{}, a tatsama: \emph{pra}, \textquotedblleft{}forth\textquotedblright{}, and \emph{sthā}, \textquotedblleft{}to stand\textquotedblright{}. CDIAL has no entry for it.
\end{cousinweb}

\begin{grammarlens}[title={What you've built}]
One more word for this chapter.
\end{grammarlens}

\subsection*{Your turn}
\begin{itemize}
\raggedright
  \item \emph{Say it:} \emph{prasthān}
  \item \emph{Say it:} \emph{prasthān}, once more
  \item \emph{Say it:} say \emph{prasthān}
  \item \emph{Recall:} say the Marathi for a supporter, then the Marathi for arrival, then say \emph{prasthān} again
\end{itemize}

\begin{tcolorbox}[breakable,colback=teal!4,colframe=teal!35!black,title={Before you move on}]
What does \mr{प्रस्थान} mean? (\textquotedblleft{}Departure\textquotedblright{}.) Say it once more.
\end{tcolorbox}
